% Curriculum vitae - Eduardo Lima Moraes
% Compile with: pdflatex cv-updated.tex
\documentclass[10pt,a4paper]{article}

\usepackage[utf8]{inputenc}
\usepackage[T1]{fontenc}
\usepackage[brazil]{babel}
\usepackage[a4paper,margin=1.55cm,top=0.75cm,bottom=1.15cm]{geometry}
\usepackage{enumitem}
\usepackage{titlesec}
\usepackage{xcolor}
\usepackage[hidelinks]{hyperref}
\usepackage{charter}
\usepackage{microtype}

\definecolor{rulegray}{HTML}{999999}
\definecolor{subgray}{HTML}{444444}

% Sections: small caps, thin rule
\titleformat{\section}{\normalsize\scshape\bfseries}{}{0pt}{}[\vspace{-7pt}\textcolor{rulegray}{\rule{\textwidth}{0.4pt}}]
\titlespacing*{\section}{0pt}{6.5pt}{3pt}

% Compact lists
\setlist[itemize]{leftmargin=1.1em,itemsep=0.7pt,parsep=0pt,topsep=1pt,label=\textbullet}

\setlength{\parindent}{0pt}
\pagestyle{empty}

% Experience/education entry
\newcommand{\entry}[4]{%
  \textbf{#1} \hfill {\small #3}\\
  {\itshape\color{subgray} #2} \hfill {\small\color{subgray}#4}\par
}

\begin{document}
\enlargethispage{3\baselineskip}

% ------------------------- HEADER -------------------------
\begin{center}
  {\huge\bfseries Eduardo Lima Moraes}\\[4pt]
  {\small
    Barueri, SP, Brasil \;\textbar\;
    \href{mailto:moraes.eduardo@proton.me}{moraes.eduardo@proton.me} \;\textbar\;
    \href{https://www.doodlesdev.com}{doodlesdev.com}\\[1pt]
    \href{https://github.com/DoodlesEpic}{github.com/DoodlesEpic} \;\textbar\;
    \href{https://www.linkedin.com/in/doodlesdev/}{linkedin.com/in/doodlesdev}
  }
\end{center}

\vspace{0pt}
% Variante compacta para uso exclusivamente digital:
% Barueri, SP, Brasil | moraes.eduardo@proton.me | doodlesdev.com | GitHub | LinkedIn

% ------------------------- SUMMARY -------------------------
\section{Resumo}
Estudante de Engenharia de Computação na Poli-USP, desenvolvendo software desde a adolescência, com experiência em sistemas embarcados, aplicações desktop e web, open source e liderança técnica. Interesse em desempenho, confiabilidade e segurança. \textbf{Busco estágio ou colaboração de meio período} em engenharia de software ou sistemas embarcados.

% Versão longa do resumo, útil para vagas mais técnicas:
% Estudante de Engenharia de Computação na Escola Politécnica da USP, desenvolvendo software desde a adolescência, com experiência prática em sistemas embarcados (STM32, FreeRTOS, C/C++), aplicações desktop multiplataforma (Qt, GTK, Rust) e web (TypeScript, Svelte). Mantenedor de projetos open source, com interesse por desempenho, confiabilidade e segurança. Busco estágio ou colaboração de meio período em engenharia de software ou sistemas embarcados.

% ------------------------- EXPERIENCE -------------------------
\section{Experiência}

\entry{Membro do Subsistema Elétrico: Sistemas Embarcados e Aquisição de Dados}{Equipe Poli Racing de Fórmula SAE}{mar.\ 2026 -- atual}{São Paulo, SP}
\begin{itemize}
  \item Lidero a recuperação e migração do datalogger para a ECU FT550 Light, atualizando firmware em C (STM32, FreeRTOS) e restabelecendo a comunicação com a ECU após descontinuidade e documentação esparsa.
  \item Desenvolvo a interface desktop em C++ (Qt 6, Qt Quick, CMake) e reescrevi o processamento de logs, reduzindo a análise de cerca de 30 minutos para 20 segundos (aprox. 90x).
  \item Atuo também no Departamento Financeiro, apoiando rotinas administrativas e o contato com capitão, diretores e patrocinadores.
\end{itemize}

\vspace{2pt}
\entry{Desenvolvedor de Software: Projeto para VichiWork}{VichiWork}{ago.\ 2019 -- fev.\ 2020}{Alphaville, SP}
\begin{itemize}
  \item Desenvolvi em equipe um protótipo web para uma empresa da área de saúde, voltado à organização e visualização de dados de planos e procedimentos, com levantamento de requisitos junto ao responsável pela empresa.
  \item Atuei na arquitetura e no frontend em Angular, implementei autenticação com Firebase Auth e configurei a aplicação no Google Cloud App Engine.
\end{itemize}

% Alternativa para vagas de software/cloud:
% \item Apoiei também a definição de infraestrutura, organização do desenvolvimento e divisão técnica das atividades da equipe, usando Git, Google Cloud e ferramentas de gestão de projeto.

\vspace{2pt}
\entry{Cofundador e Líder da Equipe de Robótica: TORC / MakerBot333}{Escola Castanheiras}{jan.\ 2019 -- nov.\ 2021}{Santana de Parnaíba, SP}
\begin{itemize}
  \item Cofundei o TORC e o Clube Maker e liderei a MakerBot333 por três anos, coordenando projetos e competições, incluindo uma sandbox de realidade aumentada e jogos com microcontroladores.
  \item Conduzi a equipe ao \textbf{5º lugar na Etapa Nacional da Olimpíada Brasileira de Robótica (OBR) de 2020}, na modalidade Simulação.
\end{itemize}

% Bullets originais do TORC, preservados para variações do currículo:
% \item Cofundei o TORC (Time Olímpico de Robótica Castanheiras) e liderei a equipe MakerBot333 durante três anos, coordenando projetos, preparação para competições e divisão de responsabilidades.
% \item Conduzi a equipe ao 5º lugar na Etapa Nacional da Olimpíada Brasileira de Robótica (OBR) de 2020, na modalidade Simulação.
% \item Cofundei o Clube Maker, desenvolvendo uma sandbox de realidade aumentada, jogos com microcontroladores para eventos escolares e outros projetos de fabricação digital.

% ------------------------- EDUCATION -------------------------
\section{Formação Acadêmica}

\entry{Bacharelado em Engenharia de Computação}{Escola Politécnica da Universidade de São Paulo (Poli-USP)}{fev.\ 2026 -- nov.\ 2030 (previsão)}{São Paulo, SP}
{\small Média ponderada: \textbf{9,0/10}. Classificação: \textbf{15º de 74} na turma de ingresso (atualizado em set.\ 2026).}\par
\vspace{2pt}
\entry{Ensino Médio}{Escola Castanheiras}{jan.\ 2019 -- nov.\ 2021}{Santana de Parnaíba, SP}
{\small Monografia: \textit{Sobre o uso da matemática na computação gráfica}. Desenvolvimento de um ray tracer educacional em C++.}\par
% Versão curta anterior: Monografia de Ensino Médio sobre renderização gráfica.
% Detalhes para variantes acadêmicas/gráficas: implementei álgebra vetorial tridimensional, câmera e composição de cena, interseções entre raios e esferas por equações quadráticas, reflexão difusa recursiva, antialiasing por amostragem estocástica e correção gama, gerando imagens no formato PPM.
% Atividades e grupos no Ensino Médio: autor e apresentador da monografia; English B Higher Level do IB Diploma Programme, nota 7; participante da 22ª Olimpíada Brasileira de Astronomia e Astronáutica (2019).

% ------------------------- PROJECTS -------------------------
\section{Projetos Open Source}
% O link explícito do GitHub no título foi removido para reduzir ruído visual; o GitHub permanece no cabeçalho.
\begin{itemize}
  \item \textbf{GraphPrime}: calculadora de números primos multiplataforma com gráficos e alto desempenho (Svelte, Tauri). \textbf{dNotes}: aplicativo de anotações multiplataforma, livre e open source (Rust, GTK 4).
  \item \textbf{Minesweeper}: jogo nativo e multiplataforma (C++, wxWidgets). \textbf{MultiplayerPicross}: nonogramas multiplayer (Svelte). \textbf{BuildUp, Billsplit, Democraz} e outros: web apps em TypeScript/Svelte, incluindo rastreador de hábitos e visualizador de projetos de lei da Câmara dos Deputados.
\end{itemize}

% Formato expandido dos projetos para vagas focadas em portfólio:
% \item \textbf{GraphPrime}: calculadora de números primos multiplataforma com gráficos e alto desempenho (Svelte, Tauri).
% \item \textbf{dNotes}: aplicativo de anotações multiplataforma, livre e open source (Rust, GTK 4).
% \item \textbf{Minesweeper}: jogo nativo e multiplataforma (C++, wxWidgets). \textbf{MultiplayerPicross}: nonogramas multiplayer (Svelte).
% \item \textbf{BuildUp, Billsplit, Democraz} e outros: web apps em TypeScript/Svelte, incluindo rastreador de hábitos e visualizador de projetos de lei da Câmara dos Deputados.

% ------------------------- SKILLS -------------------------
\section{Habilidades Técnicas}
\begin{itemize}
  \item \textbf{Linguagens:} C, C++, Rust, TypeScript/JavaScript, Python, C\#, Haskell, Shell.
  \item \textbf{Sistemas embarcados:} STM32, FreeRTOS, aquisição de dados, integração com ECU, robótica. \textbf{Desktop e web:} Qt 6 / Qt Quick, GTK 4, wxWidgets, Svelte, ASP.NET Core, Tauri.
  \item \textbf{Ferramentas:} Git, CMake, Linux. Desenvolvimento assistido por IA (agentes de código e APIs de LLMs), gestão de projetos e liderança de equipe.
\end{itemize}

% Formato expandido das habilidades:
% \item \textbf{Sistemas embarcados:} STM32, FreeRTOS, aquisição de dados, integração com ECU, robótica.
% \item \textbf{Desktop e web:} Qt 6 / Qt Quick, GTK 4, wxWidgets, Svelte, ASP.NET Core, Tauri.

% ------------------------- CERTIFICATES -------------------------
\section{Licenças e Certificados}
\begin{itemize}
  \item \textbf{Programação Funcional Pura com Aplicações}: IME-USP, curso de extensão, 30h (fev.\ 2026). Haskell.
  \item \textbf{Playgame}: Escola SAGA, 500h (2018--2019). Unreal Engine, modelagem e animação 3D e desenho.
  \item \textbf{IB English B Higher Level, nota 7}: International Baccalaureate (jul.\ 2021).
\end{itemize}

% ------------------------- LANGUAGES -------------------------
\section{Idiomas}
\textbf{Português:} nativo \quad\textbar\quad \textbf{Inglês:} avançado (IB English B HL, nota 7) \quad\textbar\quad \textbf{Espanhol:} básico a intermediário

\end{document}
